% SEGMENT OLP-0006-S01
% Part: sets-functions-relations
% Chapter: sets
% Section: subsets

\documentclass[../../../include/open-logic-section]{subfiles}

\begin{document}

\olfileid{sfr}{set}{sub}
\olsection{உட்கணங்களும் அடுக்குக் கணங்களும்}

\begin{explain}
பல சமயங்களில் கணங்களை ஒப்பிட விரும்புவோம்.
அத்தகைய ஒப்பீடுகளில் இயல்பாகத் தோன்றுவது இதுவாகும்:
\emph{ஒரு கணத்தில் உள்ள அனைத்தும் மற்றொரு கணத்திலும் உள்ளன}.
இந்த நிலைமைக்குப் போதிய முக்கியத்துவம் இருப்பதால்,
அதற்கெனப் புதிய குறியீட்டை அறிமுகப்படுத்துகிறோம்.
\end{explain}

% SEGMENT OLP-0006-S02
\begin{defn}[உட்கணம்]
$A$ என்னும் கணத்தின் ஒவ்வொரு !!{element}[um] $B$ இன்
!!a{element} எனில், $A$ என்பது $B$ இன் \emph{உட்கணம்} எனக் கூறி,
$A \subseteq B$ என எழுதுகிறோம். $A$ என்பது $B$ இன் உட்கணம் அல்ல எனில்,
$A \not\subseteq B$ என எழுதுகிறோம்.
$A \subseteq B$ ஆகவும் $A \neq B$ ஆகவும் இருப்பின்,
$A \subsetneq B$ என எழுதி, $A$ என்பது $B$ இன் \emph{தகு உட்கணம்}
எனக் கூறுகிறோம்.
\end{defn}

% SEGMENT OLP-0006-S03
\begin{ex}
ஒவ்வொரு கணமும் தனக்குத்தானே உட்கணமாகும்; $\emptyset$ என்பது
ஒவ்வொரு கணத்திற்கும் உட்கணமாகும். இரட்டை இயல் எண்களின் கணம்,
இயல் எண்களின் கணத்தின் உட்கணமாகும். மேலும்,
$\{ a, b \} \subseteq \{ a, b, c \}$.
ஆனால் $\{ a, b, e \}$ என்பது $\{ a, b, c \}$ இன் உட்கணம் அல்ல.
\end{ex}

% SEGMENT OLP-0006-S04
\begin{ex}
$2$ என்ற எண் முழுக்களின் கணத்தின் ஓர் !!{element};
இரட்டை எண்களின் கணமோ முழுக்களின் கணத்தின் ஓர் உட்கணமாகும்.
எனினும், ஒரு கணம் மற்றொரு கணத்தின் !!a{element}[pred] இருப்பதுடன் உட்கணமாகவும்
\emph{ஒரே நேரத்தில்} இருக்கக்கூடும். எடுத்துக்காட்டாக,
$\{0\} \in \{0, \{0\}\}$ என்பதும்
$\{0\} \subseteq \{0, \{0\}\}$ என்பதும் உண்மை.
\end{ex}

% SEGMENT OLP-0006-S05
கணங்கள் ஒன்றே என்பதற்கான நிபந்தனையை உறுப்புசார் சமத்துவம் தருகிறது:
$A$ இன் ஒவ்வொரு !!{element}[um] $B$ இன் !!a{element}[pred] இருப்பதுடன்,
மறுதிசையிலும் இதே நிலை அமையவும் செய்யும்போது,
அப்போதும் அப்போது மட்டுமே $A = B$.
``உட்கணம்'' என்பதன் வரையறை, இந்த நிபந்தனையின் முதல் பாதியையே
$A \subseteq B$ என வரையறுக்கிறது: $A$ இன் ஒவ்வொரு !!{element}[um]
$B$ இன் !!a{element} ஆகும்.
$A$, $B$ ஆகியவற்றை இடமாற்றினாலும் இந்த வரையறை பொருந்தும்:
அதாவது, $B$ இன் ஒவ்வொரு !!{element}[um] $A$ இன் !!a{element}[pred]
இருக்கும்போது, அப்போதும் அப்போது மட்டுமே $B \subseteq A$.
இது உறுப்புசார் சமத்துவத்தின் ``மறுதிசையிலும்'' என்ற பகுதிக்குச்
சரியாகப் பொருந்துகிறது. வேறுவிதமாகச் சொன்னால், கணங்கள் ஒன்றுக்கொன்று
உட்கணங்களாக இருக்கும்போது, அப்போதும் அப்போது மட்டுமே அவை சமம்
என்பது உறுப்புசார் சமத்துவத்திலிருந்து பின்வருகிறது.

% SEGMENT OLP-0006-S06
\begin{prop}
$A \subseteq B$, $B \subseteq A$ ஆகிய இரண்டும் நிறைவேறும்போது,
அப்போதும் அப்போது மட்டுமே $A = B$.
\end{prop}

% SEGMENT OLP-0006-S07
பயனுள்ள மேலும் சில குறியீடுகளை அறிமுகப்படுத்த இதுவும் ஏற்ற தருணமாகும்.
$A$ எப்போது $B$ இன் உட்கணமாகிறது என்பதை வரையறுக்கையில்,
``$A$ இன் ஒவ்வொரு !!{element}[um] \dots'' எனக் கூறி,
``$\dots$'' என்ற இடத்தில் ``$B$ இன் !!a{element}'' என்பதை நிரப்பினோம்.
ஆனால் இந்த வெளிப்பாட்டு \emph{வடிவம்} மிகவும் அடிக்கடி பயன்படுவதால்,
இதற்கென முறையான குறியீட்டை அறிமுகப்படுத்துவது பயனுள்ளதாக இருக்கும்.

% SEGMENT OLP-0006-S08
\begin{defn}\ollabel{forallxina}
$(\forall x \in A)\phi$ என்பது $\forall x(x \in A \lif
\phi)$ என்பதன் சுருக்கமாகும். அதேபோல், $(\exists x \in A)\phi$ என்பது
$\exists x(x \in A \land \phi)$ என்பதன் சுருக்கமாகும்.
\end{defn}

% SEGMENT OLP-0006-S09
இந்தக் குறியீட்டைப் பயன்படுத்தி, $(\forall x \in A)x \in B$
என்பது நிறைவேறும்போது, அப்போதும் அப்போது மட்டுமே $A \subseteq B$
எனக் கூறலாம்.

இப்போது ஒரு குறிப்பிட்ட வகைக் கணத்தைக் கருதுவோம்:
தரப்பட்ட கணத்தின் அனைத்து உட்கணங்களையும் கொண்ட கணமே அது.

% SEGMENT OLP-0006-S10
\begin{defn}[அடுக்குக் கணம்]
$A$ என்னும் கணத்தின் அனைத்து உட்கணங்களையும் உறுப்புகளாகக் கொண்ட கணம்
$A$ இன் \emph{அடுக்குக் கணம்} எனப்படும்; அதனை $\Pow{A}$ என எழுதுகிறோம்.
  \[
    \Pow{A} = \Setabs{B}{B \subseteq A}
  \]
\end{defn}

% SEGMENT OLP-0006-S11
\begin{ex}
$\{ a, b, c \}$ இன் சாத்தியமான உட்கணங்கள் அனைத்தும் யாவை?
அவை: $\emptyset$, $\{a \}$, $\{b\}$, $\{c\}$, $\{a, b\}$,
$\{a, c\}$, $\{b, c\}$, $\{a, b, c\}$.
இந்த உட்கணங்கள் அனைத்தையும் கொண்ட கணம் $\Pow{\{a,b,c\}}$ ஆகும்:
\[
\Pow{\{ a, b, c \}} = \{\emptyset, \{a \}, \{b\}, \{c\}, \{a, b\},
\{b, c\}, \{a, c\}, \{a, b, c\}\}
\]
\end{ex}

% SEGMENT OLP-0006-S12
\begin{prob}
$\{a, b, c, d\}$ இன் அனைத்து உட்கணங்களையும் பட்டியலிடுக.
\end{prob}

% SEGMENT OLP-0006-S13
\begin{prob}
$A$ இல் $n$ !!{element}s இருந்தால், $\Pow{A}$ இல் $2^n$
!!{element}s இருக்கும் எனக் காட்டுக.
\end{prob}

\end{document}
